\chapter{So Sánh Và Ganh Tị}
\markboth{So Sánh Và Ganh Tị}{Comparison and Envy}

\section{Người Ta Chỉ Cho Em Thấy Thành Tích, Không Cho Em Thấy Cái Giá}
\begin{truyen}{Người Ta Chỉ Cho Em Thấy Thành Tích, Không Cho Em Thấy Cái Giá}{People Show You the Achievement, Not the Price}
\chuhoa{T}{rên mạng, An thấy bạn bè khoe giấy khen, huy chương, ảnh học bổng, ảnh trường top.}
\textit{(Online, An saw friends posting certificates, medals, scholarships, and top-school photos.)}

Không ai đăng ảnh những đêm mất ngủ, những lần khóc trong nhà vệ sinh, hay những cơn hoảng loạn trước kỳ thi.
\textit{(No one posted the sleepless nights, the crying in the bathroom, or the panic attacks before exams.)}

Khi chỉ nhìn thấy phần sáng, An tưởng ai cũng đang tiến lên còn mình thì tụt lại.
\textit{(Seeing only the bright part, An believed everyone else was moving forward while she alone was falling behind.)}

Cô tư vấn học đường nói: "Con đang so hậu trường của mình với sân khấu của người khác.
\textit{(The school counselor said: "You are comparing your backstage to everyone else’s stage.)}

Cuộc so đó chắc chắn khiến con thua."
\textit{(That comparison guarantees defeat.")}

\end{truyen}

\begin{baihoc}
Khi ganh tị, hãy nhớ: em gần như không bao giờ nhìn thấy toàn bộ câu chuyện của người khác. So sánh từ dữ liệu méo thì kết luận cũng méo.
\textit{(When envy rises, remember: you almost never see the whole story of other people. If the data is distorted, the conclusion will be distorted too.)}
\end{baihoc}

\begin{ghinhoanh}
\textbf{Short English to remember:}

When envy rises, remember: you almost never see the whole story of other people.

If the data is distorted, the conclusion will be distorted too.

\end{ghinhoanh}

\begin{tuhocnhanh}
\textbf{Cau hoi tu hoc / Self-study questions}

1. Van de chinh la gi? \textit{What is the main problem?}

2. Cach xu ly tot hon la gi? \textit{What is a better action?}

3. Mot cau tieng Anh em muon nho la cau nao? \textit{Which English sentence do you want to remember?}
\end{tuhocnhanh}
\ngancach

\section{Ghen Tị Không Làm Em Xấu, Nhưng Che Nó Đi Thì Nguy}
\begin{truyen}{Ghen Tị Không Làm Em Xấu, Nhưng Che Nó Đi Thì Nguy}{Envy Does Not Make You Evil, But Hiding It Is Dangerous}
\chuhoa{K}{hi bạn cùng lớp đỗ đội tuyển, Phúc cười chúc mừng nhưng trong lòng nóng ran.}
\textit{(When his classmate got into the elite team, Phuc smiled and congratulated him, but burned inside.)}

Em ghét cảm giác đó nên cố giả như mình không quan tâm.
\textit{(He hated the feeling, so he pretended not to care.)}

Nhưng càng giả, em càng cay cú hơn.
\textit{(The more he pretended, the more bitter he became.)}

Một người thầy nói với em: "Ghen tị chỉ là tín hiệu cho thấy con đang muốn một thứ mình chưa có.
\textit{(A teacher told him: "Envy is only a signal that you want something you do not yet have.)}

Vấn đề không phải cảm xúc đó.
\textit{(The problem is not the emotion.)}

Vấn đề là con làm gì sau đó."
\textit{(The problem is what you do after it.")}

Phúc bắt đầu đổi câu hỏi từ ``Sao nó có mà mình không có?'' sang ``Nó làm được gì mà mình chưa làm?''
\textit{(Phuc began changing the question from ``Why does he have it and I do not?'' to ``What did he do that I have not yet done?'')}

\end{truyen}

\begin{baihoc}
Ghen tị được soi sáng sẽ thành động lực. Ghen tị bị giấu đi sẽ thành độc. Đừng tô hồng cảm xúc của mình; hãy đọc đúng nó.
\textit{(Envy examined becomes motivation. Envy hidden becomes poison. Do not beautify your emotions; read them accurately.)}
\end{baihoc}

\begin{ghinhoanh}
\textbf{Short English to remember:}

Envy examined becomes motivation.

Envy hidden becomes poison.

Do not beautify your emotions; read them accurately.

\end{ghinhoanh}

\begin{tuhocnhanh}
\textbf{Cau hoi tu hoc / Self-study questions}

1. Van de chinh la gi? \textit{What is the main problem?}

2. Cach xu ly tot hon la gi? \textit{What is a better action?}

3. Mot cau tieng Anh em muon nho la cau nao? \textit{Which English sentence do you want to remember?}
\end{tuhocnhanh}
\ngancach

\section{Tự Ti Không Phải Khiêm Tốn}
\begin{truyen}{Tự Ti Không Phải Khiêm Tốn}{Low Self-Worth Is Not Humility}
\chuhoa{N}{hiều học sinh nói câu này như một thói quen: ``Em dốt lắm, em không làm được đâu.''}
\textit{(Many students say this habitually: ``I am dumb. I cannot do it.'')}

Các em tưởng đó là khiêm tốn.
\textit{(They think it is humility.)}

Thực ra đó là cách đầu hàng trước khi thử thật.
\textit{(In reality it is surrender before trying seriously.)}

Khi một người tự gắn nhãn thấp cho mình, não sẽ bớt nỗ lực vì nó tin thất bại là điều đã định sẵn.
\textit{(When someone attaches a low label to the self, the brain reduces effort because failure seems already predetermined.)}

Khiêm tốn thật không phải nói mình tệ.
\textit{(Real humility is not saying you are bad.)}

Khiêm tốn thật là biết mình chưa đủ và vẫn chịu học tiếp.
\textit{(Real humility is knowing you are not enough yet and still being willing to continue learning.)}

\end{truyen}

\begin{baihoc}
Bỏ kiểu nói hạ nhục bản thân để tỏ ra an toàn. Hãy thay bằng câu trung thực hơn: ``Em chưa làm được, nhưng em có thể học.''
\textit{(Drop the self-insulting language used as a shield. Replace it with a more truthful sentence: ``I cannot do it yet, but I can learn.'')}
\end{baihoc}

\begin{ghinhoanh}
\textbf{Short English to remember:}

Drop the self-insulting language used as a shield.

Replace it with a more truthful sentence: ``I cannot do it yet, but I can learn.''

\end{ghinhoanh}

\begin{tuhocnhanh}
\textbf{Cau hoi tu hoc / Self-study questions}

1. Van de chinh la gi? \textit{What is the main problem?}

2. Cach xu ly tot hon la gi? \textit{What is a better action?}

3. Mot cau tieng Anh em muon nho la cau nao? \textit{Which English sentence do you want to remember?}
\end{tuhocnhanh}
\ngancach

\section{Đường Của Em Không Cần Giống Đường Của Bạn}
\begin{truyen}{Đường Của Em Không Cần Giống Đường Của Bạn}{Your Path Does Not Need to Match Your Friend’s Path}
\chuhoa{H}{ai người bạn thân học chung từ nhỏ.}
\textit{(Two close friends studied together from childhood.)}

Một người rất mạnh về Toán, người còn lại mạnh về Văn và giao tiếp.
\textit{(One was strong in mathematics, the other in literature and communication.)}

Đến lớp mười hai, cả hai đều cố chen vào cùng một lối vì sợ đi khác sẽ bị coi là kém hơn.
\textit{(By twelfth grade, both tried to squeeze into the same path because they feared that different meant inferior.)}

Kết quả là cả hai đều khổ.
\textit{(The result was that both suffered.)}

Người giỏi Văn học Toán trong lo âu.
\textit{(The literature student studied math anxiously.)}

Người giỏi Toán đi thi hùng biện như bị phạt.
\textit{(The math student competed in speaking as if under punishment.)}

Chỉ khi dám thừa nhận sở trường khác nhau, họ mới ngừng đánh nhau với chính thiết kế của mình.
\textit{(Only when they dared admit different strengths did they stop fighting their own design.)}

\end{truyen}

\begin{baihoc}
So sánh khiến học sinh coi khác biệt là thua kém. Nhưng khác biệt nhiều khi chỉ là dấu hiệu em nên đi con đường khác, không phải thấp hơn.
\textit{(Comparison makes students treat difference as inferiority. But often difference is only a sign that you should take a different path, not a lower one.)}
\end{baihoc}

\begin{ghinhoanh}
\textbf{Short English to remember:}

Comparison makes students treat difference as inferiority.

But often difference is only a sign that you should take a different path, not a lower one.

\end{ghinhoanh}

\begin{tuhocnhanh}
\textbf{Cau hoi tu hoc / Self-study questions}

1. Van de chinh la gi? \textit{What is the main problem?}

2. Cach xu ly tot hon la gi? \textit{What is a better action?}

3. Mot cau tieng Anh em muon nho la cau nao? \textit{Which English sentence do you want to remember?}
\end{tuhocnhanh}
\ngancach

\section{Bảng Vàng Và Cái Bóng Dưới Chân}
\begin{truyen}{Bảng Vàng Và Cái Bóng Dưới Chân}{The Honor Board and the Shadow Under It}
\chuhoa{N}{gày xưa ở làng có một bảng vàng ghi tên những người đỗ đạt.}
\textit{(Long ago in a village there was an honor board listing successful scholars.)}

Bọn trẻ lớn lên dưới cái bảng ấy, đứa nào cũng ngước nhìn.
\textit{(Children grew up under it, all of them looking upward.)}

Có đứa lấy đó làm động lực.
\textit{(Some used it as motivation.)}

Có đứa nhìn mãi rồi chỉ thấy mình thấp bé.
\textit{(Others stared so long that they felt only small.)}

Một ông đồ già nói: "Ngẩng đầu nhìn lên để biết đường là tốt.
\textit{(An old tutor said, "Looking up to know the road is good.)}

Nhưng nếu cứ vừa đi vừa ngước mãi, con sẽ vấp chân mình."
\textit{(But if you keep walking while staring upward, you will trip over your own feet.")}

So sánh chỉ hữu ích khi nó cho em phương hướng.
\textit{(Comparison helps only when it gives direction.)}

Từ lúc nó làm em ghét chính mình, nó đã hỏng.
\textit{(The moment it makes you hate yourself, it is broken.)}

\end{truyen}

\begin{baihoc}
Ngước nhìn người giỏi hơn để học là tốt. Dùng họ làm thước đo để tự chà đạp mình là dại.
\textit{(Looking at those ahead to learn is wise. Using them as a whip to beat yourself is foolish.)}
\end{baihoc}

\begin{ghinhoanh}
\textbf{Short English to remember:}

Looking at those ahead to learn is wise.

Using them as a whip to beat yourself is foolish.

\end{ghinhoanh}

\begin{tuhocnhanh}
\textbf{Cau hoi tu hoc / Self-study questions}

1. Van de chinh la gi? \textit{What is the main problem?}

2. Cach xu ly tot hon la gi? \textit{What is a better action?}

3. Mot cau tieng Anh em muon nho la cau nao? \textit{Which English sentence do you want to remember?}
\end{tuhocnhanh}
\ngancach

\section{Hai Anh Em Trong Một Nhà}
\begin{truyen}{Hai Anh Em Trong Một Nhà}{Two Siblings Under One Roof}
\chuhoa{T}{rong một nhà có hai anh em.}
\textit{(In one house lived two siblings.)}

Anh học nhanh, em chậm hơn nhưng bền.
\textit{(The older learned quickly; the younger more slowly but with endurance.)}

Người lớn khen anh trước mặt em quá nhiều.
\textit{(Adults praised the older sibling in front of the younger too often.)}

Lâu dần, em không còn nghe thấy lời khen dành cho anh nữa, chỉ nghe thấy bản án dành cho mình.
\textit{(In time, the younger no longer heard praise for the older one, only a verdict on himself.)}

Hai đứa trẻ vốn thương nhau bắt đầu khó gần nhau hơn vì mỗi lần thành công của một đứa lại thành vết xước của đứa kia.
\textit{(Two children who loved each other began drifting apart because each success of one became a scratch on the other.)}

So sánh trong gia đình không chỉ làm một đứa đau.
\textit{(Comparison in a family does not wound only one child.)}

Nó làm hỏng luôn cả tình thân.
\textit{(It damages the bond itself.)}

\end{truyen}

\begin{baihoc}
Đừng dùng một đứa trẻ để kích một đứa trẻ khác. Cách đó có thể đẩy thành tích lên chút ít nhưng làm rách mối quan hệ rất lâu.
\textit{(Do not use one child to provoke another. That method may raise performance slightly, but it tears relationships for a long time.)}
\end{baihoc}

\begin{ghinhoanh}
\textbf{Short English to remember:}

Do not use one child to provoke another.

That method may raise performance slightly, but it tears relationships for a long time.

\end{ghinhoanh}

\begin{tuhocnhanh}
\textbf{Cau hoi tu hoc / Self-study questions}

1. Van de chinh la gi? \textit{What is the main problem?}

2. Cach xu ly tot hon la gi? \textit{What is a better action?}

3. Mot cau tieng Anh em muon nho la cau nao? \textit{Which English sentence do you want to remember?}
\end{tuhocnhanh}
\ngancach

\section{Thợ Học Việc Và Người Tay Nhanh}
\begin{truyen}{Thợ Học Việc Và Người Tay Nhanh}{The Apprentice and the Quick-Handed Boy}
\chuhoa{T}{rong một xưởng mộc, một cậu học việc rất chán nản vì bạn cùng lứa cưa nhanh hơn, đục đẹp hơn.}
\textit{(In a carpentry workshop, one apprentice was discouraged because another boy his age sawed faster and carved more neatly.)}

Ông thợ cả không an ủi kiểu dễ dãi.
\textit{(The master carpenter did not offer cheap comfort.)}

Ông chỉ bảo: "Nó nhanh ở tay.
\textit{(He simply said, "He is quick with the hands.)}

Còn con chắc ở mắt.
\textit{(You are steadier with the eyes.)}

Mỗi lần đo đạc của con ít lỗi hơn nó."
\textit{(Your measurements make fewer mistakes than his.")}

Cậu bé lần đầu hiểu mình không phải bản sao lỗi của người kia.
\textit{(For the first time the boy understood that he was not a failed copy of the other.)}

Mình là một kiểu năng lực khác chưa được gọi đúng tên.
\textit{(He was a different kind of ability not yet named correctly.)}

So sánh làm mờ đi lợi thế riêng.
\textit{(Comparison blurs personal strengths.)}

Người hướng dẫn giỏi giúp học trò nhìn lại chỗ mạnh của mình.
\textit{(A good mentor helps the learner see where their own strength already lies.)}

\end{truyen}

\begin{baihoc}
Không phải ai chậm hơn ở một mặt cũng kém hơn toàn bộ. Học sinh cần được chỉ ra năng lực theo cách chính xác hơn.
\textit{(Being slower in one area does not mean being worse overall. Students need their abilities named with greater precision.)}
\end{baihoc}

\begin{ghinhoanh}
\textbf{Short English to remember:}

Being slower in one area does not mean being worse overall.

Students need their abilities named with greater precision.

\end{ghinhoanh}

\begin{tuhocnhanh}
\textbf{Cau hoi tu hoc / Self-study questions}

1. Van de chinh la gi? \textit{What is the main problem?}

2. Cach xu ly tot hon la gi? \textit{What is a better action?}

3. Mot cau tieng Anh em muon nho la cau nao? \textit{Which English sentence do you want to remember?}
\end{tuhocnhanh}
\ngancach

\section{Người Bạn Đỗ Trước}
\begin{truyen}{Người Bạn Đỗ Trước}{The Friend Who Succeeded First}
\chuhoa{H}{ai người bạn cùng ôn thi.}
\textit{(Two friends prepared for the same exam.)}

Một người đỗ sớm, một người trượt.
\textit{(One passed early; the other failed.)}

Người trượt vừa mừng cho bạn vừa đau cho mình.
\textit{(The one who failed felt happy for the friend and hurt for himself at the same time.)}

Cảm xúc đó làm em thấy tội lỗi vì nghĩ mình nhỏ nhen.
\textit{(That feeling made him guilty, as if he were petty.)}

Người bạn đỗ nói: "Mày không xấu khi thấy đau.
\textit{(The friend who passed said, "You are not bad for hurting.)}

Miễn là mày không để cơn đau đó biến thành việc ghét tao."
\textit{(As long as you do not let the pain turn into hatred toward me.")}

Nhờ câu đó, họ giữ được tình bạn và người trượt cũng giữ được lòng tự trọng của mình.
\textit{(Because of that sentence, they kept the friendship, and the one who failed also kept his self-respect.)}

\end{truyen}

\begin{baihoc}
Cảm giác chạnh lòng trước thành công của bạn là rất người. Điều quan trọng là xử lý nó tử tế.
\textit{(Feeling a sting at a friend’s success is deeply human. What matters is handling it with decency.)}
\end{baihoc}

\begin{ghinhoanh}
\textbf{Short English to remember:}

Feeling a sting at a friend’s success is deeply human.

What matters is handling it with decency.

\end{ghinhoanh}

\begin{tuhocnhanh}
\textbf{Cau hoi tu hoc / Self-study questions}

1. Van de chinh la gi? \textit{What is the main problem?}

2. Cach xu ly tot hon la gi? \textit{What is a better action?}

3. Mot cau tieng Anh em muon nho la cau nao? \textit{Which English sentence do you want to remember?}
\end{tuhocnhanh}
\ngancach

\section{Cái Giếng Làng Và Bầu Trời}
\begin{truyen}{Cái Giếng Làng Và Bầu Trời}{The Village Well and the Sky}
\chuhoa{C}{ó người chỉ quanh quẩn ở cái giếng làng nên tưởng bầu trời nhỏ lắm.}
\textit{(Some people stay around the village well and think the sky is small.)}

Có người ra khỏi làng mới biết trời rộng hơn rất nhiều.
\textit{(Those who leave discover that it is far wider.)}

Nhiều học sinh so mình với vài bạn trong lớp rồi hoặc ảo tưởng, hoặc tuyệt vọng quá sớm.
\textit{(Many students compare themselves to a few classmates and become either arrogant or hopeless far too early.)}

Thế giới lớn hơn một lớp, một trường, một kỳ thi.
\textit{(The world is larger than one class, one school, one exam.)}

Nhưng cũng vì lớn nên cuộc đua trong đầu em không nên quá chật hẹp.
\textit{(And because it is larger, the race inside your head should not be so cramped.)}

Biết mở rộng tầm nhìn giúp em bớt ganh hơn và cũng bớt tự mãn hơn.
\textit{(Broadening your view helps you become less envious and also less smug.)}

\end{truyen}

\begin{baihoc}
So sánh trong phạm vi quá hẹp dễ làm méo bản thân. Hãy mở rộng góc nhìn để thấy mình rõ hơn.
\textit{(Comparison inside a space that is too narrow easily distorts the self. Widen the view to see yourself more clearly.)}
\end{baihoc}

\begin{ghinhoanh}
\textbf{Short English to remember:}

Comparison inside a space that is too narrow easily distorts the self.

Widen the view to see yourself more clearly.

\end{ghinhoanh}

\begin{tuhocnhanh}
\textbf{Cau hoi tu hoc / Self-study questions}

1. Van de chinh la gi? \textit{What is the main problem?}

2. Cach xu ly tot hon la gi? \textit{What is a better action?}

3. Mot cau tieng Anh em muon nho la cau nao? \textit{Which English sentence do you want to remember?}
\end{tuhocnhanh}
\ngancach

\section{Người Thầy Không Xếp Hạng Tâm Hồn}
\begin{truyen}{Người Thầy Không Xếp Hạng Tâm Hồn}{The Teacher Who Refused to Rank Souls}
\chuhoa{M}{ột thầy giáo từng nói với lớp: ``Tôi phải xếp hạng bài làm của các em, nhưng tôi từ chối xếp hạng giá trị của từng đứa.''}
\textit{(A teacher once told the class, ``I may have to rank your papers, but I refuse to rank the worth of your souls.'')}

Câu nói đó làm nhiều học sinh sững lại, vì lần đầu họ nghe một người lớn tách thành tích khỏi phẩm giá.
\textit{(The sentence startled many students, because for the first time they heard an adult separate achievement from dignity.)}

Từ hôm ấy, lớp vẫn cạnh tranh nhưng bớt độc hơn.
\textit{(From that day, the class still competed, but with less poison.)}

Các em hiểu rằng thắng bài kiểm tra không cho mình quyền khinh người khác.
\textit{(They understood that winning a test did not grant the right to despise others.)}

Một môi trường học lành mạnh không cần xóa cạnh tranh.
\textit{(A healthy learning environment does not need to erase competition.)}

Nó chỉ cần chặn việc biến cạnh tranh thành tàn nhẫn.
\textit{(It only needs to stop competition from turning cruel.)}

\end{truyen}

\begin{baihoc}
Học sinh cần được nhắc rằng điểm số có thứ bậc, nhưng con người thì không nên bị đối xử như bảng xếp hạng.
\textit{(Students need to be reminded that scores have rank, but human beings should not be treated like a leaderboard.)}
\end{baihoc}

\begin{ghinhoanh}
\textbf{Short English to remember:}

Students need to be reminded that scores have rank, but human beings should not be treated like a leaderboard.

\end{ghinhoanh}

\begin{tuhocnhanh}
\textbf{Cau hoi tu hoc / Self-study questions}

1. Van de chinh la gi? \textit{What is the main problem?}

2. Cach xu ly tot hon la gi? \textit{What is a better action?}

3. Mot cau tieng Anh em muon nho la cau nao? \textit{Which English sentence do you want to remember?}
\end{tuhocnhanh}
